\documentclass[11pt]{article}
\usepackage[margin=1in]{geometry}
\usepackage{amsmath,amssymb,amsthm}
\usepackage{enumitem}
\usepackage{longtable}
\usepackage{booktabs}
\usepackage{xcolor}
\usepackage{hyperref}
\hypersetup{colorlinks=true,linkcolor=blue,urlcolor=blue,citecolor=blue}

\newcommand{\ALCIRCC}{\mathcal{ALCI}_{\mathrm{RCC5}}}
\newcommand{\EQ}{\mathrm{EQ}}
\newcommand{\DR}{\mathrm{DR}}
\newcommand{\PO}{\mathrm{PO}}
\newcommand{\PP}{\mathrm{PP}}
\newcommand{\PPI}{\mathrm{PPI}}
\newcommand{\Cl}{\operatorname{cl}}
\newcommand{\comp}{\operatorname{comp}}
\newcommand{\Reach}{\operatorname{Reach}}
\newcommand{\Mate}{\operatorname{Mate}}
\newcommand{\Mos}{\operatorname{Mos}}
\newcommand{\tp}{\operatorname{tp}}

\title{Formal Review of the Latest Repaired Split-Forest Proof Manuscript}
\author{Prepared for Michael Wessel}
\date{May 21, 2026}

\begin{document}
\maketitle

\begin{abstract}
This note reviews the latest manuscript entitled \emph{A Generated Split-Forest Decidability Proof for \(\ALCIRCC\) Without Automata: A Repaired Containment-Oriented Presentation}.  The manuscript is a significant improvement over the earlier versions: it correctly removes the semantic-Hasse/density assumption, states strong abstract RCC5 semantics, includes closure under NNF complement, separates blocking from semantic equality, and makes the multiple-upward-\(\PP\)-witness issue visible through equality-aware reachability.  Nevertheless, I do not yet regard the decidability proof as fully established.  The remaining issues are not conceptual objections to the architecture; they are formal proof gaps in the certificate semantics, the occurrence-level equality-port construction, the request-closed blocking lemma, and the support-closed mosaic/patchwork argument.  The proof is on a credible path, but several definitions still need to be tightened before the theorem can be treated as publication-grade.
\end{abstract}

\section{Overall verdict}

My current verdict is:
\[
\boxed{\text{promising repaired architecture, but not yet a complete formal decidability proof.}}
\]
The paper has moved past the earlier density/Hasse problem.  The witness-generated submodel lemma and the generated-cover reading of the containment tree are the right conceptual repairs.  The proof now resembles the originally intended split-forest proof much more closely: sparse generated witness-cover edges form the vertical backbone, split copies handle multiple containment parents, sibling information is recorded in finite interfaces, and finite certificates describe regular unfoldings.

The remaining weaknesses are concentrated in four load-bearing places:
\begin{enumerate}[label=(G\arabic*),leftmargin=*]
\item the support-closed mosaic/patchwork theorem is still too schematic;
\item equality ports and equality-aware reachability are still defined too coarsely at the state/profile level rather than fully at the occurrence/port level;
\item the request-closed cycle lemma does not yet prove actual simultaneous eventuality fulfillment in the infinite unfolding;
\item the finite certificate grammar and size bound still rely on informal ``complete profile'' and ``abstract triple graph'' notions.
\end{enumerate}
These are repairable, but they should be repaired explicitly.

\section{What is substantially improved}

The following parts of the manuscript look materially stronger than previous drafts.

\subsection{Strong abstract RCC5 semantics}

Section 2 defines strong abstract atomic RCC5 frames with
\[
\rho(x,y)=\EQ \quad\text{if and only if}\quad x=y,
\]
converse, and the RCC5 composition table.  This is the correct semantic baseline for the proof.  The displayed composition table appears consistent with the standard RCC5 base-relation composition table in the intended abstract relation-algebraic semantics.  In particular, the table now correctly has
\[
\comp(\PP,\PP)=\{\PP\},\qquad
\comp(\PPI,\PPI)=\{\PPI\},\qquad
\comp(\DR,\DR)=\{\EQ,\DR,\PO,\PP,\PPI\}.
\]

\subsection{Closure and witness thinning}

Section 3 now uses closure under NNF complement and proves the witness-generated submodel lemma in the right form.  This is an important fix.  The proof no longer depends on arbitrary semantic Hasse covers in an arbitrary starting model.  Instead, it first thins an arbitrary satisfying model to the induced submodel generated by selected witnesses for existential closure formulas.

The result is the correct bridge from arbitrary abstract satisfiability to the sparse generated model class intended by the split-forest construction.

\subsection{Generated cover edges instead of semantic Hasse edges}

The manuscript correctly states that the cover tree is not the Hasse diagram of an arbitrary semantic \(\PP\)-order.  The vertical cover edge is a generated witness-cover edge.  This removes the earlier density objection: even if an arbitrary original model has dense or non-discrete containment, the proof may pass to a witness-generated sparse submodel and then build a generated cover presentation.

\subsection{Blocking is not equality}

Section 4 explicitly separates blocking repetition from semantic equality:
\[
 u\equiv_{\mathrm{blk}}v \not\Rightarrow u\equiv_{\EQ}v.
\]
This is essential.  A finite self-loop in the certificate represents infinitely many fresh occurrence laps, not a semantic \(\PP\)-cycle and not an \(\EQ\)-collapse.

\subsection{Multiple upward witnesses are at least recognized}

The manuscript explicitly handles the known stress case of several incompatible upward \(\PP\)-witnesses by using equality-mate occurrences placed under different selected superparts.  This is the correct architectural response.

\section{Major remaining gaps}

\subsection{Gap 1: nonvertical mosaics may introduce untracked \(\PP/\PPI\) successors}

The soundness proof treats \(\PP\) and \(\PPI\) modal obligations through vertical equality-aware reachability, and treats \(\DR\) and \(\PO\) obligations through mosaics.  This division is sound only if nonvertical/sibling mosaics cannot introduce additional semantic \(\PP\) or \(\PPI\) targets that are invisible to the vertical reachability relation.

However, Definition 5.1 allows a mosaic label function
\[
L:P^2\to \{\EQ,\DR,\PO,\PP,\PPI\}
\]
with no explicit restriction that ordinary nonvertical labels must be in \(\{\DR,\PO,\EQ\}\), or that every nonvertical \(\PP/\PPI\) label is represented by the vertical reachability relation.  Then the following failure mode is possible in principle:
\begin{quote}
A mosaic supplies a \(\PP\)-label from an occurrence of state \(s\) to a side occurrence of state \(t\).  If \(\forall\PP.D\in\lambda(s)\) but \(D\notin\lambda(t)\), the semantic model produced by soundness violates the universal restriction, while clauses (V3)--(V5) may not see this target because it is not vertically reachable.
\end{quote}

The proof needs one of the following explicit repairs:
\begin{enumerate}[label=(\alph*),leftmargin=*]
\item impose a label discipline saying that every non-core sibling/frontier mosaic label is in \(\{\DR,\PO\}\), with \(\EQ\) only for explicit equality ports and \(\PP/\PPI\) only for verticalized containment links; or
\item extend \(\Reach_{\PP}\) and \(\Reach_{\PPI}\) to include all mosaic-supplied \(\PP/\PPI\) targets, and then apply (V4)--(V5) to this enlarged relation.
\end{enumerate}
The first option is closer to the original split-forest idea.

\subsection{Gap 2: equality-aware reachability is too state-level}

The certificate defines
\[
\Mate_Q(s)=\{s' : \text{some port of }s\text{ is }\equiv^Q_{\EQ}\text{-equivalent to some port of }s'\}.
\]
Soundness then assumes that for every occurrence \(u\) of state \(s\), there is a corresponding equality-mate occurrence \(\hat u\) of the required mate state \(\hat s\), so that an upward witness can be found on \(\hat u\)'s parent chain.

This is stronger than what the finite state-level definition guarantees.  A state-level equality-port link may only say that some occurrence of state \(s\) can be paired with some occurrence of state \(s'\) in some context.  It does not automatically say that every occurrence of \(s\) in every lap and branch comes with a synchronized equality companion occurrence of \(s'\).

The repair should make equality companions occurrence-local.  A certificate context should not merely record that states \(s\) and \(s'\) are in \(\Mate_Q\).  It should contain an explicit equality cluster constructor: whenever an occurrence of \(s\) is generated, the same context also generates the required mate occurrences and the explicit port links between them.  Equality-aware reachability should be port- and context-indexed, for example
\[
(s,p)\;\Reach_R\;(t,q),
\]
rather than just state-indexed.  The soundness theorem can then legitimately say that the mate occurrence exists in the unfolding.

\subsection{Gap 3: the request-closed repetition lemma is not strong enough}

Lemma 8.3 is meant to replace a Buchi-style acceptance argument.  This is acceptable in principle, but the current formulation is too weak.  It says that a request may be either fulfilled inside the segment or passed out at the same profile/status, and then argues that repetition gives a fresh opportunity to fulfill it.

A fresh opportunity is not the same as fulfillment.  Consider a cycle carrying a request \(\exists\PP.D\) but containing no strictly later position with \(D\).  Repeating the cycle forever does not satisfy the request; it merely postpones it forever.

The cycle condition should be strengthened to an actual fulfillment condition.  For every cycle occurrence/state/port carrying a transitive request \(\exists R.D\), the finite cyclic segment must contain a strict cyclic successor, possibly in the next lap, reachable through the allowed equality-aware mechanism, whose type contains \(D\).  Formally, for each generated request at a cycle node \(c_i\), there should be some \(j>i\) modulo the cycle length, with at least one strict step, such that
\[
 c_i\;\Reach_R^{+}\;c_j
 \quad\text{and}\quad
 D\in\lambda(c_j).
\]
If equality mates are used, the source of the strict path must be an explicit mate occurrence in the same equality cluster.

The current graph reachability clause (V4) points in this direction, but Lemma 8.3 and its proof should be rewritten so that request-closed cycles mean actual cyclic fulfillment, not merely repeated deferral.

\subsection{Gap 4: the support-closed mosaic/patchwork theorem is underdefined}

Definitions 5.1--5.2 and Lemma 5.3 are conceptually right: mosaics must be joint support objects, not raw pair tables.  However, the proof still hides a large amount of work inside phrases such as ``the abstract triple graph generated by those symbols is finite'' and ``by support closure, this abstract triple occurs in a legal mosaic.''

For the theorem to be fully checkable, the manuscript should define:
\begin{enumerate}[label=(\alph*),leftmargin=*]
\item the finite set of abstract occurrence shapes produced by the unfolding scheme;
\item the quotient map from every concrete occurrence triple in the infinite unfolding to one abstract triple shape;
\item the label assignment function induced by mosaics and its overlap-agreement condition;
\item the exact finite support-closure test over those abstract shapes.
\end{enumerate}
Then Lemma 5.3 can be proved by a simple lifting argument: every concrete violating triple maps to a finite abstract violating triple, which is rejected by the finite test.

As written, SC4 is close to being circular: it requires that every abstract triple that can be produced is represented, but the abstract triple universe itself is not formally constructed.

\subsection{Gap 5: DR/PO witness placement remains schematic}

Lemma 4.3 states that selected \(\DR\) and \(\PO\) witnesses are placed inside finite side interfaces or sibling mosaics.  Clause (V6) later requires a named concrete mosaic position witnessing each \(\DR/\PO\) existential.  This is the right condition.

What remains missing is a construction showing that every selected \(\DR/\PO\) witness from the witness-generated model can be placed into the finite side-interface system while preserving all interactions with containment ancestors, descendants, equality mates, and other side witnesses.  The manuscript says that labels are inherited from the original model, so each finite side network is legal.  That is plausible, but the finite construction still needs to explain which boundary nodes are included and why the chosen width suffices uniformly.

\subsection{Gap 6: the finite bound is still not fully derived}

Lemma 9.1 gives a computable bound \(B(C_0)\), but the proof relies on a context width that is described only informally:
\begin{quote}
``one source, one witness position for each existential formula, one equality-mate boundary position for each containment witness mechanism, and the positions needed for all triples of such positions.''
\end{quote}
This is not yet a precise bound because the number of equality-mate boundary positions and the number of mosaic positions depend on the not-yet-formal support-closure and equality-cluster constructions.

This is unlikely to be a fatal problem: some very large primitive-recursive or even naive enumeration bound should suffice.  But the proof should give a definite finite grammar first, then define \(B(C_0)\) as the number of grammar objects up to the chosen arities.  At present the bound depends on ingredients that have not been fully specified.

\section{Smaller technical issues}

\begin{enumerate}[label=(S\arabic*),leftmargin=*]
\item \textbf{Identity should be included in mate reachability.}  \(\Mate_Q(s)\) should explicitly include the source occurrence itself, or the identity port of the source, otherwise ordinary vertical witnesses with no equality splitting may accidentally require an explicit port.

\item \textbf{State-level strict reachability may be too coarse for equality exclusion.} Clause (V9)(Eb) forbids equality between states connected by strict \(\PP/\PPI\)-reachability.  In cyclic or branching finite quotients, state-level reachability may identify profiles that occur both comparably and noncomparably in different concrete places.  The exclusion should be stated at the level of occurrence ports in contexts, not only states.

\item \textbf{Support closure should distinguish ordinary sibling nodes from core/equality nodes.}  The intended \(\DR/\PO\) sibling simplification applies to ordinary non-core frontier nodes after containment and equality copies are accounted for.  The current mosaic definition should reflect this stratification explicitly.

\item \textbf{Soundness induction should explicitly cover all semantic successors.}  In the \(\PP/\PPI\) universal case, the proof should state and prove that every semantic \(\PP/\PPI\)-successor in the quotient arises from equality-aware vertical reachability.  Without the nonvertical-label discipline from Gap 1, this statement is false.

\item \textbf{Completeness should separate extraction from regularization.}  The proof of Lemma 8.4 mixes extraction of representative contexts, finite blocking, equality-port recording, and mosaic collection.  It would be clearer and safer to split this into separate lemmas: finite profile extraction, lasso/blocking, equality-port preservation, and mosaic support closure.
\end{enumerate}

\section{Concrete recommended repairs}

I recommend the following concrete changes to the manuscript.

\subsection{Repair R1: add a verticalization discipline}

Add a validity clause such as:
\begin{quote}
Every occurrence pair labelled \(\PP\) or \(\PPI\) in the unfolded weak structure is either a strict vertical ancestor/descendant pair, or is obtained from such a pair by replacing one or both endpoints by explicit equality mates.  Ordinary sibling-frontier mosaic labels are restricted to \(\DR\) or \(\PO\).
\end{quote}
Then prove the lemma:
\[
 [u]R[v],\; R\in\{\PP,\PPI\}
 \quad\Longrightarrow\quad
 \text{some equality mate of }u\text{ is vertically }R\text{-reachable to a representative of }[v].
\]
This lemma is exactly what the soundness induction needs.

\subsection{Repair R2: make equality clusters explicit constructors}

Replace the state-level \(\Mate_Q(s)\) relation by explicit finite equality clusters attached to context occurrences.  The finite certificate should say:
\begin{quote}
Whenever an occurrence of source state \(s\) is generated, the context simultaneously generates mate positions \((s_1,p_1),\ldots,(s_k,p_k)\) with explicit \(\EQ\)-ports connecting them, and all external labels are congruent across the cluster.
\end{quote}
Then equality-aware reachability becomes a relation from a source occurrence-port to a target occurrence-port.  This will remove the current gap between ``some state can be a mate'' and ``this occurrence has that mate available.''

\subsection{Repair R3: strengthen request-closed cycles to fulfilled cycles}

Define a finite cycle to be request-accepting only if every transitive existential request generated at a cycle position is actually satisfied at a strict later cyclic position, perhaps through an explicit equality mate.  Do not allow indefinitely repeated deferral unless the repeated lap itself contains the desired target type.

For a cycle \(c_0,\ldots,c_{m-1}\), a sufficient check is:
\[
\forall i\;\forall(\exists R.D\in\lambda(c_i))\;
\exists j\;\bigl(c_i\Reach_R^+ c_j\text{ in the cyclic unfolding and }D\in\lambda(c_j)\bigr).
\]
This is finite and does not require naming Buchi automata.

\subsection{Repair R4: define the finite abstract triple graph}

Before stating support closure, define the finite graph \(\mathcal T_Q\) whose vertices are abstract triple shapes generated from:
\begin{itemize}[leftmargin=*]
\item profile names;
\item context names;
\item port names;
\item relative vertical positions: same occurrence, ancestor, descendant, sibling interface, equality mate;
\item cycle-offset class where needed for strict cyclic reachability.
\end{itemize}
Then require \(\Mos_Q\) to contain a legal representative for every triple shape in \(\mathcal T_Q\), with restriction maps agreeing on overlaps.  Lemma 5.3 will then become a finite abstraction argument rather than an informal patchwork claim.

\subsection{Repair R5: give a grammar-first decidability proof}

Define the finite certificate grammar first, including maximum arities and port counts.  Then the decidability proof can simply enumerate all grammar objects up to that bound.  Avoid relying on a ``complete boundary profile'' whose components are described semantically before their finite syntax is fixed.

\section{Revised status assessment}

The manuscript is now a coherent and serious proof proposal.  The following obstacles from earlier versions appear resolved or substantially resolved:
\begin{center}
\begin{tabular}{@{}ll@{}}
\toprule
Earlier issue & Current status \\
\midrule
Semantic Hasse/density objection & Resolved by witness thinning and generated covers \\
Weak closure for universal duals & Resolved by NNF-complement closure \\
Strong equality not stated & Resolved in the frame definition \\
Blocking treated as equality & Conceptually resolved \\
Multiple upward witnesses ignored & Recognized and architecturally addressed \\
Incorrect RCC5 table & Corrected \\
\bottomrule
\end{tabular}
\end{center}

But the following remain open at proof level:
\begin{center}
\begin{tabular}{@{}ll@{}}
\toprule
Remaining issue & Severity \\
\midrule
Nonvertical \(\PP/\PPI\) labels invisible to reachability & High \\
Occurrence-level existence of equality mates & High \\
Actual fulfillment of cyclic requests & High \\
Formal abstract triple graph for mosaics & High \\
Uniform finite side-interface width & Medium-high \\
Precise enumeration bound & Medium \\
\bottomrule
\end{tabular}
\end{center}

\section{Conclusion}

I would not yet state that the manuscript establishes decidability.  I would state instead:
\begin{quote}
The latest manuscript gives a credible repaired split-forest proof architecture for decidability of \(\ALCIRCC\) over strong abstract RCC5 frames.  The main conceptual objections from earlier drafts have been addressed.  However, the proof still requires additional formalization of occurrence-level equality clusters, verticalization of all \(\PP/\PPI\) semantic targets, fulfilled-cycle blocking, and the finite support-closed mosaic/patchwork construction.  Once those items are made precise, the proof could plausibly become a complete decidability proof without resorting to explicit automata machinery.
\end{quote}

The strongest next step is not another high-level rewrite, but a targeted tightening of the definitions around contexts, ports, cycles, and abstract triple shapes.  Those are the places where the current proof still risks accepting a finite certificate whose unfolding does not satisfy the intended semantic conditions.

\end{document}
